\documentclass[11pt,letterpaper]{article}
\usepackage[T1]{fontenc}
\usepackage{lmodern,amsmath,amssymb,amsthm,mathtools}
\usepackage[margin=1in]{geometry}
\usepackage{microtype,booktabs,longtable,array,enumitem,xcolor}
\usepackage{xurl}
\usepackage[colorlinks=true,linkcolor=blue!45!black,citecolor=blue!45!black,urlcolor=blue!45!black]{hyperref}
\usepackage{fancyhdr}
\pagestyle{fancy}\fancyhf{}
\fancyhead[L]{\small Asymmetric packet balancing for subexponential SAT}
\fancyhead[R]{\small Research memorandum}
\fancyfoot[C]{\thepage}
\setlength{\headheight}{14pt}
\setlength{\emergencystretch}{2em}
\newtheorem{theorem}{Theorem}[section]
\newtheorem{lemma}[theorem]{Lemma}
\newtheorem{proposition}[theorem]{Proposition}
\newtheorem{corollary}[theorem]{Corollary}
\theoremstyle{remark}\newtheorem{remark}[theorem]{Remark}
\newcommand{\Ztwo}{\mathbb Z_{(2)}}
\newcommand{\NN}{\mathbb Z_{\ge0}}
\newcommand{\poly}{\operatorname{poly}}
\newcommand{\polylog}{\operatorname{polylog}}
\newcommand{\vTwo}{\nu_2}
\newcommand{\ceil}[1]{\left\lceil #1\right\rceil}
\newcommand{\floor}[1]{\left\lfloor #1\right\rfloor}
\newcommand{\sTwo}{s_2}
\newcommand{\KL}{K_{\mathrm{asym}}}
\newcommand{\sourceurl}[1]{\url{#1}}
\hypersetup{pdftitle={Asymmetric packet balancing for subexponential SAT},pdfauthor={Research analysis prepared for Vladimir Reshetnikov},pdfsubject={Review of seven ETH proposals, deterministic balancing, Cauchy modes, explicit contraction bounds}}
\title{Asymmetric packet balancing for\newline subexponential SAT\\[4pt]\large A review of the recent ProveIt proposals and a sharper contraction bound}
\author{Research analysis prepared for Vladimir Reshetnikov\\\small with ChatGPT}
\date{October 26, 2026 (Pacific time)}
\begin{document}
\maketitle
\begin{abstract}
We review the seven analysis packets placed in ProveIt's
\texttt{Computability/SubexponentialSAT} project on October 26, 2026.
Their strongest justified general-density bound remains
$\poly_k(n,m)2^{O_k(n\log\log n/\sqrt{\log n})}$, with a sharper
clause-density dependence and several improvements to explicit coefficients.
We prove a further refinement of their common tensor-contraction framework.
A deterministic, exact-rational, two-stage labeling replaces collision
variance by nearly uniform pair loads. A second improvement uses the
asymmetric degrees created by the recursive child construction. Together
with an additive precision buffer, the previously proved sharp Hermite
budget, and factorial Cauchy modes, the resulting normalized bound is
\[
 T(s,\mathcal L)\le(s+\mathcal L)^{O(1)}
 2^{(\KL+o(1))s/\sqrt{\log_2s}},\qquad
 \KL=\frac{\pi\sqrt{5/2}}{\ln2}(9+\sqrt3)
 =76.9089949516\ldots .
\]
Here the root has at most $s$ Boolean wires and maximum degree at most
$\lceil s^{9/10}\rceil$, with the other input conditions of the cited
contraction theorem. This improves the earlier explicit certificate
$648.4410666488\ldots$ in the same size normalization; the comparison
concerns upper bounds, not measured runtimes. A separate integral Cauchy
basis exploiting the actual reciprocal row domain gives mode slopes
$4,16/5,64/21,\ldots$, without increasing the number of mode families.
We also give a $31.0641320630\ldots$ coefficient for suitably low-degree
inputs. These results do not improve the power of the logarithm or remove
the general-density $\log\log n$ factor. The report contains ordinary
mathematical proofs and exact finite prototypes. It does not supply an
end-to-end SAT solver or a new proof-assistant certificate, and its SAT
consequences retain the stated upstream dependencies.
\end{abstract}
\noindent\textbf{Source boundary.} The preprint is pinned to
\texttt{openai/math@7ef5c0cc5c}. The repository placement and comparison
are pinned to \texttt{ProveIt@2b925d8800}. Full identifiers, immutable links,
and file hashes are supplied in \texttt{source\_manifest.json}.
\newpage
\tableofcontents
\newpage

\section{Outcome, comparison, and meaning of the improvement}
\label{sec:outcome}
The requested target was an improvement beyond the recent ETH-related
proposals, not merely a restatement of the cited preprint. The new
contribution is a different constructive packing lemma and a recursion
schedule that uses the two sides' different degree bounds. It composes
with the strongest already proposed algebraic refinements.

The following theorem is the principal result. All logarithms in its
normalization are base two; $\ln$ is the natural logarithm.
\begin{theorem}[Refined contraction bound, relative to the stated framework]
\label{thm:main}
Let $s\ge2$. Consider a bipartite Boolean tensor network with at most $s$
wires, polynomially many vertices, maximum degree at most
$\Delta=\lceil s^{9/10}\rceil$, and total graph/tensor description length
$\mathcal L$. Its entries belong to
\[
 (\mathbb Z/2^{P_0}\mathbb Z)[V]/(V^\ell),\qquad
 1\le P_0,\ell\le C s,
\]
for a fixed $C$. Suppose each original tensor entry is evaluable in
$(s+\mathcal L)^b2^{A s^{9/10}}$ time and $(s+\mathcal L)^b$ space,
for fixed $A,b$. Using the exact deformation and child-contraction
identities of the cited preprint, there is a deterministic algorithm
with polynomial space and bit time
\begin{equation}\label{eq:headline}
 (s+\mathcal L)^{O(1)}
 2^{(\KL+o(1))s/\sqrt{\log_2s}},\qquad
 \KL=\frac{\pi\sqrt{5/2}}{\ln2}(9+\sqrt3).
\end{equation}
The $o(1)$ term tends to zero for the fixed interface constants above.
The algorithm may use either the earlier factorial Cauchy basis or the
stronger annulus basis proved in Section~\ref{sec:annulus-modes}.
\end{theorem}

The qualifier in the theorem is substantive. We prove the new packing,
precision scheduling, and complexity statements here; we retain the
preprint's exact algebraic contraction architecture. We inspected the
relevant source identities and the earlier modifications, but did not
rebuild the entire upstream formalization or independently formalize
this theorem.

\subsection{What the constants compare}
The contraction size $s$ counts Boolean wires. It is not automatically
the number $n$ of SAT variables. For a general formula, the reduction
has $s=O_k(n\log\log n)$, with a width-dependent constant that must be
tracked before converting any numerical wire coefficient into a
numerical coefficient on $n\log\log n$.

\begin{center}
\begin{tabular}{@{}p{.60\textwidth}r@{}}
\toprule
Normalized contraction certificate & Base-two coefficient\\
\midrule
Earlier source-schedule analysis in the contraction packet & $18535.7538\ldots$\\
Earlier refined schedule \cite{contractionpacket} & $648.4411\ldots$\\
Earlier schedule plus factorial modes, matched composition & $374.4$\\
New symmetric balancing and additive-buffer schedule & $93.9291\ldots$\\
New asymmetric recursive schedule, Theorem~\ref{thm:main} & $76.9090\ldots$\\
New low-degree specialization, Appendix~\ref{app:lowdegree} & $31.0641\ldots$\\
\bottomrule
\end{tabular}
\end{center}
The low-degree row has stronger input assumptions and is listed separately.
The numerical comparisons use the common normalized interface with
$P_0,\ell\le s$; the new theorem additionally permits any fixed
$P_0,\ell\le Cs$.
The matched composition in the third row uses the earlier refined
schedule, whose leaf cost is already lower order, and multiplies its
mode coefficient by $1/\sqrt3$. Its unrounded value is
$374.378\ldots$ and is recorded by the verification script. It is a composition for comparison, not another
historical packet. The final coefficient is about $11.86\%$ of the
previous $648.44$ certificate. This is not an $8.43$-fold wall-clock
speedup: an exponent coefficient is being compared, and upper bounds do
not establish ratios of actual running times.

\subsection{What remains unresolved}
For fixed clause width $k$, the strongest justified density envelope
among the reviewed packets is still
\begin{equation}\label{eq:density-envelope}
 \poly_k(n,m)\,
 2^{O_k\left(\frac{n}{\sqrt{\log n}}
      \min\{\log(2+m/n),\log\log n\}\right)}.
\end{equation}
Our engine can be used inside that reduction and inside the scalar parity
lift of the earlier refinement. It changes their contraction constants,
not the asymptotic expression in~\eqref{eq:density-envelope}.
The proposed $n/\log n$ and related stronger exponents in the packets
rest on explicitly missing summation or factorization interfaces. We do
not treat those conditional targets as established algorithms, nor claim
to have constructed their missing interfaces.

\section{The reviewed commits and proposal frontier}
\label{sec:review}
The initially retrieved snapshot was
\texttt{aa6933442eb936962de79ba6d321cd0391c450bd}. During the review,
commit \texttt{2b925d8800cda68f5723b65dc1fcdaf3d8170359}
placed the same delivered material under
\texttt{Computability/SubexponentialSAT/Research}. We reviewed that
placement and its index as well as the original archives. The arrival
commits were:
\begin{itemize}[leftmargin=1.5em]
\item \texttt{a505bd768bce347c70e923586c5f72ad3085bf44}: three packets
and the two primary manuscript PDFs;
\item \texttt{64b3202ee84fc4d71fe69c34d2110c93b7da95cd}: three further
refinement packets;
\item \texttt{8afdbbea48d67b8d252a99e8926f3eca9649eae3}: the asymptotic
analysis and barrier packet.
\end{itemize}
The placement records that these are external research artifacts, with
no ProveIt Lean or Rocq development and no executed package validation
at placement time \cite{repoindex}. Our checks below are new, scoped
checks of our own components; they do not change that historical status.

\begin{longtable}{@{}>{\raggedright\arraybackslash}p{.24\textwidth}>{\raggedright\arraybackslash}p{.43\textwidth}>{\raggedright\arraybackslash}p{.25\textwidth}@{}}
\toprule
Packet & Strongest proved refinement claimed in the packet & Remaining limitation\\
\midrule\endfirsthead
\toprule Packet & Strongest proved refinement claimed in the packet & Remaining limitation\\\midrule\endhead
Factorial Cauchy modes \cite{cauchypacket} & Weights
$1+b+2\nu_2(b!)$; leading abstract mode entropy decreases by
$1/\sqrt3$. & Same logarithmic order; original packing retained.\\[5pt]
Contraction improvements \cite{contractionpacket} & Leading precision
$6MD/u$; explicit coefficient $648.441\ldots$; streamed packet
coefficient extraction. & Same global logarithmic order; packet factor
three retained.\\[5pt]
Precision refinement \cite{precisionpacket} & Exact Hermite condition
count and loss; real taper minimax calculation. & Stronger structured-mode
summation is an unproved hypothesis.\\[5pt]
Uniform density improvement \cite{densitypacket} &
Equation~\eqref{eq:density-envelope}, uniformly in clause density. &
No stronger unrestricted $n$-only exponent.\\[5pt]
Density refinement \cite{densityrefinement} &
$n\log(2+B)/\sqrt{\log n}$ for $m\le Bn$ in its stated range;
precision/schedule savings. & Overlaps the preceding density result;
no new logarithmic power.\\[5pt]
Scalar parity refinement \cite{scalarpacket} & Scalar integral lift;
incidence-degree cap $k(p+a-1)$ rather than the source's
$k(a+2(p-1))$. & Recursive compression requires an unsupplied compact
factorization.\\[5pt]
Asymptotic analysis \cite{barrierpacket} & Scoped obstructions and precise
missing interfaces; no stronger unrestricted bound claimed. & Neither
linear-wire compression nor degree-preserving recursion is constructed.\\
\bottomrule
\end{longtable}

All seven analyze the same four preprint source files: the two pinned
upstream commits in their manifests have identical relevant blobs.
The separate 49-page manuscript credited to Claude (Anthropic),
\emph{A $2^{o(n)}$ time algorithm for 3-SAT}, states a bound with exponent
$O(\min\{km/\sqrt{\log m},k^2n\log\log n/\sqrt{\log n}\})$
\cite{claudepaper}. It is a different construction. We inspected its
headline and scope; we do not claim an explicit constant comparison to
that manuscript or certify its reported unpublished Lean development.

The OpenAI manuscript's sharper appendix gives the general exponent
$n\log\log n/\sqrt{\log n}$ and fixed-density exponent
$n/\sqrt{\log n}$ \cite{upstream,reduction}. The upstream formalization
scope page states a weaker $(\log\log n)^2/\sqrt{\log n}$ numerator
factor in its general bound \cite{leanscope}. Improving that weaker
formal endpoint by quoting the manuscript would not answer the present
research question.

\section{The inherited exact contraction interface}
\label{sec:interface}
A Boolean tensor network has one Boolean variable per edge, with
contraction
\[
 Z=\sum_{x\in\{0,1\}^{E}}\prod_v T_v(x|_{\delta(v)}).
\]
Parallel edges have distinct ports. The source works over an unramified
extension $\mathcal O$ of $\Ztwo$ whose residue field has polynomial size,
and over $\mathcal A=\mathcal O[V]/(V^\ell)$ modulo powers of two.
Valuation and divisibility are coordinatewise in the fixed free bases.
The arithmetic cost is polynomial in all the displayed parameters.

At an active call let $M$ be its wire count, $L=2^a$ its packet length,
and $r$ its fixed alphabet size. Labels belong to $[r]^a$ on both sides;
write $Q=r^a$. Every ordered pair of labels is padded to the same load
$H_{\rm pad}$, a multiple of $L$. The source's orientation construction
then groups the padded slots into packets of length $L$. In coordinate
$i$, the unordered pair of symbols is fixed within a packet, and bit $i$
of a slot determines its orientation when the two symbols differ.
Unused slots have no real ports.

The deformation uses the integer grid
\[
 S=\{-L,-L+1,\ldots,L-1\},\qquad
 \mathcal C(T)=\prod_{t\in S}(T-t),\qquad \alpha=2^{a+1},
\]
and a packet wire kernel
\begin{equation}\label{eq:cauchy-interface}
 W_{xy}(T)=\frac{\mathcal C(T)}{\mathcal C'(0)(T-(y-x))}.
\end{equation}
The denominator's linear factor cancels a numerator factor.
The deformed contraction $F(T)$ has degree at most $M(2L-1)$,
$F(0)$ is the desired contraction, and each $F(t+\alpha z)$ is integral.
At $i=\nu_2(t)$, the specialization $W(t)$ connects only ports in the
same coordinate-$i$ group.

The nonlocal terms at $T=t+\alpha\zeta$ separate into linear source
moments and target coefficients. At level $i$ a column tile has step
$2^{i+1}$ and length $J=D_i/2$, where $D_i=L/2^i$. The earlier factorial
basis has certified weights
\begin{equation}\label{eq:factorial-weights}
 w_j=1+j+2\nu_2(j!)=3j+1-2s_2(j),\qquad j\ge0.
\end{equation}
A tuple $\beta$ contributes only when $\sum_Zw_Z\beta_Z<p_t$, where
$p_t$ is the working precision. For each tuple the source's exact
factorization is
\begin{equation}\label{eq:recursion-interface}
 F(t+\alpha\zeta)=\sum_{\beta:\sum w_Z\beta_Z<p_t}
                  \prod_{g=1}^r Z_{g,\beta}\pmod{2^{p_t}}.
\end{equation}
The child for group $g$ puts \emph{all retained tensors from both sides}
of the parent on one side, and new packet tensors on the other.
A new packet tensor has degree at most $2L$.

This last structural fact drives the new schedule. Below the root, only
one side can contain the original high-degree tensors. The other side is
always made of fresh, much smaller packet tensors.

\begin{remark}[Canonical representatives]
Each tensor description retains its creation precision. Evaluating an
old tensor later means evaluating it at that original precision and
reusing its canonical integer coordinates. Canonical lifting is not a
ring homomorphism. Neither the packing change nor the new mode basis
licenses silently changing a tensor's defining lift. Our consequences
use exactly the source's immutable-description convention.
\end{remark}

\section{Deterministic asymmetric balancing of pair loads}
\label{sec:packing}
We now prove the central new combinatorial lemma independently of tensor
algebra. It applies to multigraphs and uses exact rational arithmetic.

\begin{theorem}[Two-stage pair-load balancing]\label{thm:pair-balance}
Let a bipartite multigraph have $M$ edges, maximum row degree at most
$R$, and maximum column degree at most $C$, where $R,C$ are positive
integers. For integers $Q\ge1$
and $h\ge2$, there is a deterministic labeling of each side by $Q$ labels
such that every ordered label-pair load $X_{ab}$ satisfies
\begin{equation}\label{eq:pair-load}
 X_{ab}\le\left(1+\frac{2}{h-1}\right)
              \left(\frac{M}{Q^2}+\frac{R}{Q}\right)
          +2(hC+1)\ln Q.
\end{equation}
The construction runs in polynomial bit time and space in the explicit
graph size and the parameters $Q,h,R,C$.
\end{theorem}
\begin{proof}
First assign row vertices greedily to the label with smallest current
sum of incident degrees. If $W_a$ is the resulting stub load of row
label $a$, then
\[
 W_a\le M/Q+R.
\]
Indeed, just before the final vertex placed in a bin, that bin has at
most the current average load, and the vertex adds at most $R$.

For each column vertex $v$, let $w_{av}$ count its edges to row label
$a$. Multiplicities are counted, so $0\le w_{av}\le C$ and
$\sum_vw_{av}=W_a$. Give each column vertex an independent uniform
random label temporarily. Put
\[
 q=1+\frac1{hC},\qquad \Phi=\sum_{a,b=1}^Q q^{X_{ab}}.
\]
For a fixed pair,
\[
 \mathbb E q^{X_{ab}}
   =\prod_v\left(1+\frac{q^{w_{av}}-1}{Q}\right).
\]
Convexity of $q^w$ on $[0,C]$ gives
$q^w-1\le(w/C)(q^C-1)$. Hence
\begin{equation}\label{eq:expected-potential}
 \mathbb E\Phi\le Q^2
 \exp\left(\frac{(q^C-1)\max_aW_a}{QC}\right).
\end{equation}

Assign the column labels in a fixed order. For each choice compute the
conditional expectation of $\Phi$ exactly and choose a minimum. The
average over the $Q$ choices equals the old conditional expectation, so
the selected choice does not increase it. At termination,
$q^{\max X_{ab}}\le\Phi\le\mathbb E\Phi$. Taking logarithms for the
analysis, not for the algorithm, yields
\[
 \max X_{ab}\le\frac{2\ln Q}{\ln q}
       +\frac{q^C-1}{C\ln q}\frac{\max_aW_a}{Q}.
\]
Now $1/\ln q\le hC+1$, since $\ln(1+x)\ge x/(1+x)$.
Also $q^C\le e^{1/h}$, and $e^{1/h}-1\le1/(h-1)$, so
\[
 \frac{q^C-1}{C\ln q}
 \le\frac{h+1/C}{h-1}\le1+\frac2{h-1}.
\]
This proves~\eqref{eq:pair-load}, including $Q=1$.

For the algorithmic claim, a conditional expectation is a sum of terms
$q^{x_{ab}}\prod_{v\ \mathrm{unassigned}}\frac{Q-1+q^{w_{av}}}{Q}$.
All entries are rational, and the total exponent of $q$ in any product
is at most $M$. Products over polynomially many vertices therefore have
polynomial bit length. There are $Q$ candidate labels and $Q^2$ pair
terms at every step. Straightforward exact numerator/denominator
arithmetic already gives polynomial bit complexity; no numerical
exponential or logarithm is computed by the labeling algorithm.
\end{proof}

\begin{corollary}[Nearly lossless uniform padding]\label{cor:padding}
Let $h=\max\{2,\lceil\log_2s\rceil\}$ and $Q=s^{O(1)}$.
Suppose, uniformly over the active calls under consideration,
\begin{equation}\label{eq:padding-conditions}
 \frac{QR}{M}=o(1),\qquad
 \frac{Q^2C(\log s)^2}{M}=o(1),\qquad
 \frac{LQ^2}{M}=o(1).
\end{equation}
Round the largest pair load upward to a multiple $H_{\rm pad}$ of $L$.
Then
\begin{equation}\label{eq:padded-mass}
 P:=H_{\rm pad}Q^2=(1+o(1))M.
\end{equation}
\end{corollary}
\begin{proof}
Multiply~\eqref{eq:pair-load} by $Q^2$ and divide by $M$.
Each excess term tends to zero by~\eqref{eq:padding-conditions}; rounding
adds less than $LQ^2$. Conversely $P\ge M$, because it contains all real
edges. This proves the equality with a uniform relative error.
\end{proof}

\subsection{Why the power-scale conditions are useful and close to sharp}
With $R=C=\Delta$, the earlier second-moment proof required
$Q^4\Delta=o(M)$ for nearly lossless padding. The new lemma needs only
$Q^2\Delta(\log s)^2=o(M)$, in addition to harmless rounding and the
weaker row term. Thus its admissible leading packet exponent is twice
as large in the symmetric setting.

There is also a scoped obstruction. A row vertex of degree $R_0$ forces
some pair load to be at least $R_0/Q$, regardless of the labels, so
uniform padding uses at least $QR_0$ slots. A parallel bundle of
$\mu$ edges forces at least $Q^2\mu$ padded slots. Consequently
$QR_0=O(M)$ and $Q^2\mu=O(M)$ are necessary in worst-case examples for
this particular common-pair-load architecture. Taking $\mu=C$ when
$R\ge C$ witnesses the column-degree scale. Our sufficient conditions
match these powers with logarithmic overhead and strict asymptotic
slack. The little-oh conditions themselves are not claimed necessary,
and this is not a lower bound for tensor contraction or SAT.

\section{Packet counts, crossing families, and the recursive asymmetry}
\label{sec:packets}
The exact orientation construction can be retained after replacing its
label assignment. The common padded mass $P=H_{\rm pad}Q^2$ gives
$P/L$ packets. In each coordinate and on each side, a fixed symbol has
exactly $P/r$ padded ports. Therefore its real-port count is at most
$P/r$, and a child has at most
\begin{equation}\label{eq:child-size}
 2P/r=(2/r+o(1))M
\end{equation}
wires.

Only packets with two different coordinate symbols create crossing
modes. Exactly a fraction $1-1/r$ of the padded ordered label pairs
have this property. Each packet is homogeneous for it, so at valuation
level $i$ the number of tile families is at most
\begin{equation}\label{eq:family-count}
 c_i\le(1-1/r)\frac{P}{L}\,2^{i+1}
     =\left(2(1-1/r)+o(1)\right)\frac{M}{D_i}.
\end{equation}
This retains the exact equal-symbol saving as well as the improvement
from padded mass $3M$ to $(1+o(1))M$.

\begin{lemma}[Side-degree invariant]\label{lem:side-degree}
At depth $d\ge1$, orient the child bipartition with all retained tensors
on the row side and all fresh packet tensors on the column side.
If all tensors created so far have degree at most $\Delta$, then
\[
 R_d\le\Delta,\qquad C_d\le2L_{d-1}.
\]
These bounds hold regardless of the chosen mode tuple, interpolation
node, or interpolation value.
\end{lemma}
\begin{proof}
The source child construction places all retained parent tensors from
both old sides onto one new side. None gains a port. Each vertex on the
other side is newly created for a parent packet and has at most its
$L_{d-1}$ row and $L_{d-1}$ column ports. The degree is independent of the
coefficient values and of the mode indices. The same construction repeats
at every subsequent call.
\end{proof}

\section{Precision and entropy with an additive buffer}
\label{sec:precision}
The previous packets already sharpened the Hermite precision from a
leading $24MD/u$ to $6MD/u$ \cite{contractionpacket,precisionpacket}.
We use that result, with its proof ingredients explicit here, and change
how successive packet exponents accommodate inherited precision.

At valuation $i<u\le a$, let $q_i=\rho(u-i)$, with $q_i=0$ for $i\ge u$.
There are $L/2^i$ nonzero grid points of valuation $i$. Thus the exact
condition count is
\begin{equation}\label{eq:condition-count}
 \sum_{t\ne0}q_t=\rho L(2u-2+2^{1-u}).
\end{equation}
Choose
\begin{equation}\label{eq:taper}
 \rho=\left\lceil
 \frac{M(2L-1)+1}{L(2u-2+2^{1-u})}
 \right\rceil.
\end{equation}
The denominator is an integer because $u\le a$. The count exceeds the
degree of $F$, and when $u\to\infty$, $u/M\to0$,
$\rho=(1+o(1))M/u$ uniformly.

The interpolation multiplier at a point of valuation $i$ has negative
valuation
\[
 \Lambda_i=\rho\sum_{j=1}^{a-i}
 j\min\{j,u-i\}\,2^{a-i-j}\le6\rho D_i,
\]
using $\sum_{j\ge1}j^2/2^j=6$. The inverse local Hermite factor is
integral, and evaluating a jet of degree below $q_i$ loses at most
$(q_i-1)(a+1-i)$ more bits. Hence the simple sufficient precision
\begin{equation}\label{eq:working-precision}
 p_t=p+6\rho D_i+\rho u(a+1)
\end{equation}
preserves the required residue modulo $2^p$. The exact smaller loss can
also be used; the theorem needs only this upper envelope.

Let $\lambda=\lceil\log_2s\rceil$ and
\begin{equation}\label{eq:additive-buffer}
 b=\lceil6\log_2\lambda\rceil,
 \qquad a_d=a_{d-1}+u_d+b,\qquad a_{-1}=0.
\end{equation}
Assume active sizes are $M>s/\lambda$, the depth is $O(\log\lambda)$,
and $a_d,u_d=O(\lambda)$ with $u_d\to\infty$ uniformly. For large $s$,
$\rho\le2M/u$. Every precision increment is therefore bounded by
$O(s2^{a_j}+s\lambda)$ at its ancestor depth $j$. Summing gives the
explicit, sufficient estimate
\begin{equation}\label{eq:inherited}
 p=O(s\lambda^2 2^{a_{d-1}}).
\end{equation}
At the root it follows from $P_0=O(s)$.
Since $i<u_d$,
\[
 D_i\ge2^{a_d-u_d+1}=2^{a_{d-1}+b+1}.
\]
Combining with $s/M\le\lambda$, $u_d\le\lambda$ gives
\[
 \frac{p}{MD_i/u_d}=O(\lambda^4 2^{-b})=O(\lambda^{-2})=o(1).
\]
Also $\rho u(a+1)/(MD_i/u)=O(u(a+1)/D_i)=o(1)$.
Consequently
\begin{equation}\label{eq:sharp-precision}
 p_t\le(6+o(1))\frac{MD_i}{u_d}
\end{equation}
uniformly. The old cumulative rule doubled every retained range;
only an additive $O(\log\log s)$ buffer is required.

\subsection{The factorial weighted-partition bound}
For the weights in~\eqref{eq:factorial-weights},
\[
 \sum_{j\ge0}-\ln(1-e^{-\theta w_j})
   =\frac{\pi^2+o(1)}{18\theta},\qquad\theta\downarrow0.
\]
Indeed $w_j=3j+O(\log(j+2))$; compare it outside a fixed initial segment
with $(3\pm\varepsilon)j$, use
$\int_0^\infty-\ln(1-e^{-x})\,dx=\pi^2/6$, and then send
$\varepsilon$ to zero. The finite initial segment contributes
$O(\log(1/\theta))$.
For at most $c_i$ families, the number of tuples of weighted degree less
than $p_t$ is bounded by
\[
 N_i\le\exp\left(\theta p_t+
             \frac{(\pi^2+o(1))c_i}{18\theta}\right).
\]
Extending a finite family to all $j\ge0$ only increases this bound.
Optimize $\theta$ using the upper envelopes for $p_t$ and $c_i$ in
\eqref{eq:family-count} and~\eqref{eq:sharp-precision}. The chosen
$\theta$ tends uniformly to zero, since it is proportional to
$\sqrt{u_d}/D_i$ and $D_i\ge\lambda^6$. We obtain
\begin{equation}\label{eq:local-cost}
 \ln N_i\le(\kappa_r+o(1))\frac{M}{\sqrt{u_d}},\qquad
 \kappa_r=\pi\sqrt{\frac83(1-1/r)}.
\end{equation}
No lower bound on the actual family count is needed: its prescribed
upper envelope can be used directly in the exponential estimate.

\section{The explicit asymmetric schedule and its global bound}
\label{sec:asymmetric-schedule}
For Theorem~\ref{thm:main}, set
\begin{align}
 r&=16,& H&=\lceil s/\lambda\rceil,& b&=\lceil6\log_2\lambda\rceil,\notag\\
 u_0&=\lfloor\lambda/81\rfloor,&
 u_d&=\left\lfloor\frac{\lambda}{108\,4^{d-1}}\right\rfloor\ (d\ge1),&
 a_d&=\sum_{j=0}^d(u_j+b).
 \label{eq:main-schedule}
\end{align}
The sum of the leading retained ranges after the root equals
$\lambda/81$, so the cumulative budget is $2\lambda/81$.

\subsection{Applicability without a circular depth assumption}
First consider every hypothetical depth
$d\le d_*:=\lceil\log_6\lambda\rceil+2$ and every active $M>H$.
Uniformly in that range,
\[
 a_0\le\lambda/81+O(\log\lambda),\quad
 a_d\le2\lambda/81+O((\log\lambda)^2),\quad
 u_d\ge\lambda^{\eta+o(1)},\quad
 \eta=1-\frac{\ln4}{\ln6}>0.
\]
Thus flooring is a uniform relative $o(1)$ perturbation.

At the root, both degree bounds are $\Delta$. Since $Q=16^{a_0}$,
\[
 Q^2\Delta=s^{809/810+o(1)},
 \qquad \frac8{81}+\frac9{10}=\frac{809}{810}<1.
\]
The other terms in~\eqref{eq:padding-conditions} are smaller. Below the
root, Lemma~\ref{lem:side-degree} gives $R_d\le\Delta$ and
$C_d\le2^{a_{d-1}+1}$. Their relevant leading exponents are
\begin{center}
\begin{tabular}{@{}lll@{}}
\toprule
Quantity & Upper power of $s$, up to $o(1)$ & Exact margin\\
\midrule
$Q_dR_d$ & $4(2/81)+9/10=809/810$ & $1/810$\\
$Q_d^2C_d$ & $8(2/81)+2/81=2/9$ & $7/9$\\
$L_dQ_d^2$ & $9(2/81)=2/9$ & $7/9$\\
\bottomrule
\end{tabular}
\end{center}
These fixed gaps absorb the logarithmic factors, the additive buffers,
and the factor $s/M\le\lambda$. New packet degrees are below $\Delta$
because $2/81<9/10$. Corollary~\ref{cor:padding} therefore holds
uniformly for every hypothetical active depth under consideration.

For sufficiently large $s$, each child has at most $M/6$ wires by
\eqref{eq:child-size}. An active path cannot then exceed
$\lceil\log_6\lambda\rceil+O(1)$. This validates the hypothetical range
and closes the induction. More sharply, for a uniform $\delta_s\to0$,
\begin{equation}\label{eq:sharp-shrink}
 M_d\le s\left(\frac{1+\delta_s}{8}\right)^d.
\end{equation}
All exceptional bounded sizes, or sizes at which the explicit parameter
checks fail, can be handled by direct summation. The theorem is
asymptotic; the strict but small root gap does not suggest practical
input thresholds.

\subsection{Summing the actual work}
For $r=16$, \eqref{eq:local-cost} gives
$\kappa_{16}=\pi\sqrt{5/2}$. The root contributes
$(9+o(1))s/\sqrt\lambda$ to $\sum_dM_d/\sqrt{u_d}$.
The tail contributes at most
\[
 (1+o(1))\frac{s\sqrt{108}}{\sqrt\lambda}
 \sum_{d\ge1}\frac{2^{d-1}}{8^d}
 =\left(\sqrt3+o(1)\right)\frac{s}{\sqrt\lambda}.
\]
The perturbation in~\eqref{eq:sharp-shrink} changes a strictly convergent
geometric series by $o(1)$. Thus
\begin{equation}\label{eq:path-sum}
 \sum_d\frac{M_d}{\sqrt{u_d}}
 \le(9+\sqrt3+o(1))\frac{s}{\sqrt\lambda}.
\end{equation}
The polynomial number of interpolation nodes, values, and group choices
at each of $O(\log\lambda)$ depths adds only
$O(\lambda\log\lambda)$ to the logarithm of the total work.
Direct leaves cost $2^{O(s/\lambda)}$; original entry routines cost
$\exp(O(s^{9/10}))$; packet entry routines cost
$\exp(O(L\log s))=\exp(s^{2/81+o(1)})$.
These are all $2^{o(s/\sqrt\lambda)}$.
Multiplying the per-call costs over the branching tree and dividing the
natural logarithm by $\ln2$ proves~\eqref{eq:headline}.

The total label sets and interpolation degrees are polynomial in $s$;
in particular $M(2L-1)<s^4$ eventually. Precisions remain polynomial by
\eqref{eq:inherited} and $a_d=O(\lambda)$. Old descriptions are immutable
and shared by reference. Enumerate tuples, interpolation values, and
child groups depth first, accumulating one sum and product at a time.
This preserves polynomial working space and finishes the proof of
Theorem~\ref{thm:main}.

% Original research note for integration into the parent report.
% This file is an input fragment, not a standalone preamble.
\section{An additional improvement from the actual Cauchy row domain}
\label{sec:annulus-modes}

The factorial Cauchy packet in the reviewed proposals uses the enlarged
row domain $z\in2\mathbb Z_{(2)}$ (up to a highly divisible extension-ring
perturbation). The actual reciprocal rows exclude a residue class.
Keeping this information gives a termwise stronger integral mode basis
without adding any source-tile families. The result is local: its worst
limiting slope over all valuation levels remains three.

Write $\nu_2(0)=+\infty$, and let
\[
 S_i=\{w\in\mathbb Z_{(2)}:0\le\nu_2(w)\le i\},\qquad i\ge0.
\]
Define $G(b)=b+\nu_2(b!)$ for $b\ge0$. The next construction is explicit;
no theorem about generalized factorials is needed.

\begin{lemma}[An explicit ordering of a punctured two-adic ball]
\label{lem:annulus-order}
For each $i\ge0$ there are distinct positive integers
$a^{(i)}_0,a^{(i)}_1,\ldots$ in $S_i$ and nondecreasing integers
$F_i(0),F_i(1),\ldots$ such that
\[
 F_i(b)=\min_{w\in S_i}\sum_{h<b}\nu_2(w-a^{(i)}_h),
 \qquad
 \sum_{h<b}\nu_2(a^{(i)}_b-a^{(i)}_h)=F_i(b).
\]
They can be constructed as follows. For $i=0$ take
$a^{(0)}_b=2b+1$ and $F_0(b)=G(b)$. For $i\ge1$, the sequence
$F_i$ is the sorted merge, with multiplicity, of the two strictly
increasing sequences
\[
 G(n),\qquad n+F_{i-1}(n),\qquad n=0,1,2,\ldots.
\]
When the next element comes from the first sequence output the point
$2n+1$; when it comes from the second output $2a^{(i-1)}_n$.
Resolve a tie in favor of the first sequence. For every $i,b$,
\[
 a^{(i)}_b\le2b+1,
 \qquad F_i(b)\ge\nu_2(b!).
\]
\end{lemma}

\begin{proof}
For the ordinary consecutive ordering $0,1,\ldots$, the minimum
valuation of a product of $n$ consecutive differences over
$\mathbb Z_{(2)}$ is $\nu_2(n!)$. This follows by counting, for every
$r\ge1$, the at least $\lfloor n/2^r\rfloor$ factors in the appropriate
residue class; equality holds at the next consecutive integer.
Scaling and translating gives the assertion for odd integers, including
rational two-integral arguments.

Split $S_i$ into its odd part $1+2\mathbb Z_{(2)}$ and its even part
$2S_{i-1}$. Differences between the two parts are units. If $n_o$
odd and $n_e$ even points have already been chosen, the least possible
valuation for the next odd point is $G(n_o)$, and that for the next even
point is $n_e+F_{i-1}(n_e)$. Choosing the smaller proves the asserted
minimum and the merge recursion. Both component sequences are strictly
increasing: $G$ increases by at least one, as does $n+F_{i-1}(n)$.

For the lower bound, use induction on $i$. The second component is
termwise at least $G$ once $F_{i-1}(n)\ge\nu_2(n!)$. The sorted merge
of two copies of $G$ is precisely $\nu_2(b!)$, because
$\nu_2((2n)!)=\nu_2((2n+1)!)=G(n)$. Thus the merge in the statement
is no smaller. The base case is immediate.

The same termwise comparison and the odd-first tie rule imply that
before an even point is selected, $n_o\ge n_e+1$. For an odd point,
$2n_o+1\le2b+1$. For an even point, induction bounds it by
$2(2n_e+1)\le2(n_o+n_e)=2b$. This proves the point-size bound and
also makes the construction polynomial in the number of requested
points and $i$.
\end{proof}

\begin{lemma}[Uniform valuation slope]
\label{lem:annulus-slope}
Put
\[
 R_i=1+\frac{3}{4^{i+1}-1}.
\]
Then, for all $i,b\ge0$,
\[
 R_i b-4\log_2(b+2)\le F_i(b)\le R_i b.
\]
\end{lemma}

\begin{proof}
The base case follows from
$F_0(b)=2b-s_2(b)$, where $s_2(b)$ is the binary digit sum.
For the induction step put $B=1+R_{i-1}$, so $2<B\le3$ and
\[
 R_i=\left(\frac12+\frac1B\right)^{-1}
     =\frac{2B}{B+2}.
\]
Substitution gives the displayed closed form for $R_i$.
The two component sequences are bounded above by $2n$ and $Bn$.
Their merge at index $b$ is at most $R_i b$: the number of terms in
these two arithmetic sequences no greater than $R_i b$ is at least
$b+1$, since the two unrounded ratios sum to the integer $b$.

For the lower bound let $x=F_i(b)$ and
$L_b=\log_2(b+2)$. Each of the first $b+1$ terms of the merge has a
component index at most $b$. The lower bounds for those components
are $2n-L_b$ and $Bn-4L_b$. Consequently
\[
 b+1\le \frac{x+L_b}{2}+1+\frac{x+4L_b}{B}+1,
\]
so
\[
 x\ge R_i b-R_i-\left(\frac{R_i}{2}
                      +\frac{4R_i}{B}\right)L_b.
\]
Here $R_i\le6/5$ and
$R_i/2+4R_i/B=(B+8)/(B+2)\le5/2$.
Since $L_b\ge1$, the last two losses together are at most
$4L_b$. This proves the uniform estimate.
\end{proof}

\begin{lemma}[Stability under an unramified extension perturbation]
\label{lem:annulus-extension}
Let $\mathcal O$ be the unramified coefficient ring of the contraction
argument, with valuation extending $\nu_2$. Suppose $R\ge2$,
$w_0\in S_i$, $w-w_0\in2^R\mathcal O$, and
$0\le b<2^{R-2}$. Then
\[
 \nu_2\left(\prod_{h<b}(w-a^{(i)}_h)\right)\ge F_i(b).
\]
\end{lemma}

\begin{proof}
The assertion is immediate for $b=0$. Otherwise the selected integers
are distinct and satisfy $1\le a^{(i)}_h\le2b-1<2^{R-1}$.
Choose the ordinary residue $c\in\{0,\ldots,2^R-1\}$ of $w_0$.
If $c$ is nonzero and is not a selected point, take $w'=c$.
It belongs to $S_i$, and all factor valuations are below $R$ and
unchanged by the perturbation. If $c=0$, then $\nu_2(w_0)\ge R$,
so $i\ge R$; take $w'=2^R\in S_i$. Again all differences to the
selected points have valuation below $R$.
Finally, if $c=a^{(i)}_h$ for one selected point, take
$w'=a^{(i)}_h+2^R\in S_i$. Its difference from that point has
valuation $R$, and its other difference valuations agree with those
of $w$; the corresponding difference for $w$ has valuation at least
$R$. In each case the product valuation for $w$ is at least that for
$w'\in S_i$, which is at least $F_i(b)$ by
Lemma~\ref{lem:annulus-order}. This argument uses only ordinary
residues of $w_0$; $w$ itself may have nonzero extension coordinates.
\end{proof}

\begin{theorem}[Annulus Cauchy modes]
\label{thm:annulus-modes}
At a contraction node of valuation $i$, retain the existing packet
layout and column tiles. Put
\[
 \alpha=2^{a+1},\quad T=t+\alpha\zeta,\quad
 R=a+1-i,\quad J=2^{a-i-1},\quad y=y_0+2^{i+1}j,
\]
where $0\le j<J$ and $\zeta\in\mathcal O$. For a row crossing
that tile, let
\[
 d_x=T-y_0+x,\quad z_x=\frac{2^{i+1}}{d_x},\quad
 A_x=\frac{\mathcal C(T)}{\mathcal C'(0)d_x},
 \qquad \xi_b=2a^{(i)}_b.
\]
Define
\[
 r_b(z)=\frac{b!\prod_{h<b}(z-\xi_h)}
                    {\prod_{h=0}^{b}(1-hz)},
 \qquad
 s_b(j)=\frac{(1-b\xi_b)\binom{j}{b}}
                    {\prod_{h=0}^{b}(1-\xi_h j)}.
\]
The exact finite identity
\[
 W_{xy}(T)=A_x\sum_{b=0}^{J-1}r_b(z_x)s_b(j)
\]
holds. All source functions $s_b(j)$ are two-integral; all displayed
new denominators are units. The target coefficient $A_xr_b(z_x)$ is
divisible by $2^{\omega_i(b)}$, where
\[
 \boxed{\ \omega_i(b)=R+b+F_i(b)+\nu_2(b!)\ }.
\]
In particular these weights are termwise no smaller than the previous
factorial weights with the same base contribution $R$, and
\[
 \omega_i(b)=R+\left(3+\frac{3}{4^{i+1}-1}\right)b
                 +O(\log(b+2)),
\]
with an absolute error constant. For $i=0$ the exact formula is
$\omega_0(b)=R+2b+2\nu_2(b!)$.
\end{theorem}

\begin{proof}
The crossing condition gives $e=\nu_2(t-y_0+x)\le i$. Since
$\alpha\zeta$ has valuation at least $a+1$, one has
$\nu_2(d_x)=e$, so $z_x$ is divisible by two and
$\nu_2(A_x)\ge a+1-e\ge R$, as in the original integrality proof.
For
\[
 w=z_x/2,\qquad w_0=2^i/(t-y_0+x),
\]
one has $w_0\in S_i$ and
\[
 \nu_2(w-w_0)\ge i+a+1-2e\ge a+1-i=R.
\]
Lemma~\ref{lem:annulus-extension} therefore yields
\[
 \nu_2\left(\prod_{h<b}(z_x-\xi_h)\right)
 \ge b+F_i(b)\qquad (b<J=2^{R-2}).
\]
The remaining denominator factors in $r_b$ are units, proving the
weight bound after multiplying by $b!$ and $A_x$.

For completeness the identity is a telescoping Cauchy elimination.
Let
\[
 U_b(z)=\prod_{h<b}\frac{z-\xi_h}{1-hz},\qquad
 V_b(j)=\prod_{h<b}\frac{j-h}{1-\xi_hj}.
\]
Then direct multiplication gives
\[
 \frac{U_b(z)V_b(j)-U_{b+1}(z)V_{b+1}(j)}{1-zj}
 =\frac{(1-b\xi_b)U_b(z)V_b(j)}{(1-bz)(1-\xi_bj)}
 =r_b(z)s_b(j).
\]
For $0\le j<J$ the factor $V_J(j)$ vanishes. Summation gives
$1/(1-zj)=\sum_{b<J}r_b(z)s_b(j)$, which becomes the kernel
identity after multiplication by $A_x$. Since every $\xi_b$ is even,
all source denominators and $1-b\xi_b$ are units. The binomial
coefficient is evaluated as an ordinary integer: no factorial is
divided in the residue ring. The comparison with the earlier weights
uses $F_i(b)\ge\nu_2(b!)$, and the slope follows from
Lemma~\ref{lem:annulus-slope} and
$\nu_2(b!)=b-s_2(b)$.
\end{proof}

\paragraph{Algorithmic consequence and its scope.}
Use one mode for each existing packet, tile, and index $b<J$, with
source linear form $\sum_j s_b(j)X_{y_0+2^{i+1}j}$. Its source and
target groups are unchanged. Hence the number of families, child
network sizes, entry-evaluation mechanism, and canonical-lift
invariant are unchanged. Truncate by the new exact weights.
The admissible tuple set is contained in that certified by the old
factorial weights, for every finite parameter choice. Precomputing
the ordering and evaluating the integral features costs only a
polynomial factor in the existing parameters. The uniform entropy
analysis may still use slope three, while a fixed valuation level
has the sharper slope
$\gamma_i=3+3/(4^{i+1}-1)$. These slopes approach three as $i$ grows.
Thus this result by itself is not a new logarithmic power in the
worst-case SAT exponent, nor does it prove a uniform leading constant
strictly better than the slope-three envelope over all levels.
A valuation-dependent taper may exploit the local improvements, but
that optimization is not assumed in the global theorem here.


\section{Consequences for SAT and integration with the earlier results}
\label{sec:sat}
For bounded clause density, the source uses a linear-size bounded-degree
network. A clause contributes its satisfaction tensor. A variable's
occurrences are linked by a tree of equality tensors, and an additional
subdivision makes the network bipartite. Every satisfying assignment has
one consistent extension to the internal equality wires. Unused variables
contribute a power of two. Contracting modulo $2^{n+1}$ recovers the exact
integer count because it lies in $[0,2^n]$.

The new packing and schedule preserve this encoding and exactness.
Moreover the low-degree theorem in Appendix~\ref{app:lowdegree} applies
to the original bounded-degree networks. The guaranteed time is still
$\poly_k(n,m)2^{O_{k,D}(n/\sqrt{\log n})}$ when $m\le Dn$.
The numerical coefficient in terms of variables additionally depends
on the wire count of the chosen encoding; the wire-level constants in
this report must not be substituted directly for it.

For arbitrary density, the earlier reductions use sparse parity tests
and an incidence subset representation with $O_k(n\log\log n)$ wires.
The uniform density-aware packet improves that count in its stated
regimes, yielding~\eqref{eq:density-envelope}. The scalar parity packet
reduces the retained incidence-degree cap and removes the auxiliary
coefficient axis for that lift. These changes are compatible with our
engine: it permits $\ell=1$ as well as truncated extension rings, and its
inputs are still arbitrary tensors meeting the stated oracle bounds.
Their correctness and density constructions remain separate imported
results. This report adds no new sparsification or expander-code theorem.

For integration, replace the preprint's collision labeling by
Theorem~\ref{thm:pair-balance}, retaining the same packet orientation.
Keep the exact side provenance of every child; do not discard it into a
single maximum-degree parameter. Use~\eqref{eq:main-schedule}, the taper
\eqref{eq:taper}, and the precision rule~\eqref{eq:working-precision}.
Replace the mode basis by either~\eqref{eq:factorial-weights} or the
annulus basis. The earlier streamed coefficient-extraction routine is
also compatible, provided its temporary extra precision does not change
the tensor's creation-time representatives. A separate
\texttt{INTEGRATION.txt} gives the source-label map.

\section{Executable evidence and its limits}
\label{sec:checks}
The deliverable contains original standard-library Python components for
pair-label balancing, packet reconstruction, annulus orderings and exact
Cauchy identities, and schedule arithmetic. The recorded run summary is
included as \texttt{results/verification\_summary.json}. The tests use
exact integers, rational numbers, and modular arithmetic for mathematical
identities. Decimal constants are displays of exact formulas, not inputs
to a proof or to the labeling algorithm.

The packing prototype computes every candidate conditional expectation
as a rational, checks the averaging identity, chooses a minimum, and
checks the final potential. Multiedges and asymmetric degree fixtures
are included. An independent reconstruction counts edges by their
returned labels; a packet checker verifies real-edge preservation,
shared slots, coordinate orientations, group load bounds, and the
unequal-packet fraction on small examples.

The annulus checks exercise selected-point minimality, the merge
recurrence, the point-size and slope bounds, rational Cauchy identities,
and arithmetic over
\[
 (\mathbb Z/2^{160}\mathbb Z)[U]/(U^2+U+1).
\]
The nontrivial extension coordinates are essential: arbitrary extension
residues do not satisfy ordinary integer-valued binomial divisibility.
The proof uses the actual row's highly divisible congruence to a rational
point in the permitted annulus.

For a concrete comparison of the new modes, use one tile family with
$J=16$, base valuation $R=6$, and precision $p=64$. Compare with the
previous factorial weights using the \emph{same} base contribution $R$.
Exact dynamic programming gives:
\begin{center}
\begin{tabular}{@{}crr@{}}
\toprule
Valuation level $i$ & Previous feasible tuples & Annulus feasible tuples\\
\midrule
0 & 1726 & 634\\
1 & 1726 & 1421\\
2 & 1726 & 1683\\
3 & 1726 & 1723\\
\bottomrule
\end{tabular}
\end{center}
These are counts in a specified abstract mode-budget model, not running
times or numbers of necessarily nonzero SAT contributions. The choice
$a=i+5$ makes the displayed $J$ and $R$ consistent with the packet geometry.

The schedule checker verifies the exact rational exponent margins,
budget identities, convergence conditions, and numerical evaluations
of the displayed constants. Its finite parameter examples are labeled
as diagnostics; many small logarithmic sizes fail the asymptotic
conditions and are not represented as successful practical instances.

These are component checks and proof-supporting fixtures. The all-size
conclusions come from the proofs. There is no full SAT engine in the
archive, no SAT performance benchmark, no compiled Lean or Rocq theorem,
and no verification of the entire upstream dependency graph.

\section{Why the logarithmic order survives, and further research}
\label{sec:questions}
The weighted enumeration still has linearly growing mode weights.
With $c\asymp M/D$, $p\asymp MD/u$, and $w_j\asymp j$, choose
$B$ to be a sufficiently small fixed multiple of $D/\sqrt u$.
There are $\Theta(cB)=\Theta(M/\sqrt u)$ mode variables in indices
$[B,2B)$. Every zero-one choice among them fits the budget after adjusting
the fixed multiple. Thus the abstract mode pool can have
$\exp(\Omega(M/\sqrt u))$ feasible tuples. Their total unweighted degree
is $O(M/\sqrt u)\le M$, so imposing the affine-product degree cap alone
does not remove the obstruction. This is a scoped enumeration example,
not a statement that every such tuple is nonzero on every actual SAT
network, and not a lower bound for all algorithms.

The annulus slopes tend to three as $i$ grows. Consequently a uniform
replacement of three by a larger fixed slope has not been obtained.
The new packing conditions still constrain $u$ to $O(\log s)$ in this
architecture. Together these observations explain why the present
result changes the leading coefficient while preserving the square-root
logarithmic denominator.

The most concrete next questions are:
\begin{enumerate}[leftmargin=1.6em]
\item \textbf{Can feasible modes be summed without enumerating them?}
A useful theorem must provide an actual uniform factorization or
coefficient algorithm for the packet-group tensors, with bounded wire
count and controlled oracle complexity. Merely assigning shorter names
to the tuples is insufficient.
\item \textbf{Can the dense incidence representation lose its extra logarithm?}
An $O(n)$-wire representation, with valid degree and entry-cost bounds,
would remove the current $\log\log n$ overhead. The earlier packets state
this missing interface; the present work does not construct it.
\item \textbf{Can instance-sensitive degree and multiplicity information
improve the first call?}
The root alone pays the symmetric high-degree constraint. A construction
that uses smaller parallel multiplicity, local neighborhood balance, or
large-degree special structure could improve this remaining coefficient
bottleneck. It would need a polynomial-time construction, not just a
random-existence estimate.
\item \textbf{What is the best finite taper with annulus weights?}
The exact sequences $F_i(j)$ give more information than their limiting
slopes. A finite optimization that incorporates actual family counts,
node losses, and inherited precision could reduce enumeration in the
specified finite model. It must not be advertised as a new logarithmic
order unless a uniform asymptotic argument establishes that.
\item \textbf{Can the new pieces be formalized and integrated?}
The rational conditional-expectation lemma, packet invariants, and finite
Cauchy identity are distinct formalization targets. A full result also
needs the cost model, exact quotient-ring arithmetic, inherited
precision, and immutable tensor descriptions. The supplied prototypes
are a starting specification, not substitutes for those proofs.
\end{enumerate}

\appendix
\section{A low-degree contraction coefficient}
\label{app:lowdegree}
When the root has maximum degree at most $\lceil s^{1/12}\rceil$,
one can instead maintain that same small degree bound throughout and
use symmetric packing. Set
\[
 r=48,\qquad u_d=\left\lfloor\frac{\lambda}{14\,8^d}\right\rfloor,
 \qquad a_d=\sum_{j\le d}(u_j+b),\quad
 b=\lceil6\log_2\lambda\rceil,\quad H=\lceil s/\lambda\rceil.
\]
Then $a_d\le4\lambda/49+O((\log\lambda)^2)$.
The new packet degree is at most $s^{4/49+o(1)}<s^{1/12}$.
Since $3^5<2^8$, one has $\log_2 48<28/5$, and
\[
 \frac8{49}\log_2 48+\frac1{12}
 <\frac8{49}\frac{28}{5}+\frac1{12}
 =\frac{419}{420}<1.
\]
Thus $Q^2\Delta\polylog(s)=o(M)$ and $LQ^2=o(M)$ uniformly at active
calls. Coarse shrinkage $M/12$ and hypothetical depth
$\lceil\log_{12}\lambda\rceil+2$ give the same noncircular bootstrap;
$u_d$ grows at least as a fixed positive power of $\lambda$ there.
The sharp child ratio is $(1+o(1))/24$, the family coefficient is
$47/24$, and the natural-log mode coefficient is
$\pi\sqrt{47/18}$. Therefore
\begin{equation}\label{eq:lowdegree-constant}
 K_{\rm low}=\frac{\pi\sqrt{47/18}\sqrt{14}}
 {\ln2(1-\sqrt8/24)}=31.06413206299\ldots .
\end{equation}
Original entry routines may still have the source's allowed
$\exp(O(s^{9/10}))$ cost: that remains lower order. Bounded-degree
clause networks satisfy the smaller graph-degree premise eventually.
The lower coefficient is not asserted for arbitrary high-degree inputs.

For a direct intermediate comparison under the generic $s^{9/10}$
degree premise, replace $14$ in this schedule by $128$. Then
$a_d\le\lambda/112+O((\log\lambda)^2)$ and
$2\log_2(48)/112+9/10<1$. The coefficient becomes
$K_{\rm sym}=93.9291064374\ldots$, the symmetric row of the comparison
table. The main theorem gains further by using the fresh side's smaller
degree instead of imposing this generic bound on both sides.

\section{Optimization of the depth allocation}
\label{app:optimization}
For a fixed child ratio $c\in(0,1)$ and budget
$\sum_du_d\le A\lambda$, the continuous minimization problem is
\[
 \min\sum_{d\ge0}\frac{c^d}{\sqrt{u_d}}.
\]
H\"older's inequality, applied to
$c^{2d/3}=(c^d/\sqrt{u_d})^{2/3}u_d^{1/3}$, gives
\[
 \sum_d\frac{c^d}{\sqrt{u_d}}
 \ge\frac{1}{\sqrt{A\lambda}(1-c^{2/3})^{3/2}}.
\]
Equality is attained by
$u_d=A\lambda(1-c^{2/3})c^{2d/3}$.
This optimizes a stated schedule model, not all contraction algorithms.

For the asymmetric high-degree model, the root budget is
$A_0=(1-\gamma)/(2\log_2r)$ and the retained-side cumulative budget
below the root is $2A_0$, provided the fresh-side and rounding constraints
also hold. When $c^{2/3}<1/2$, the root cap is active. Spend $A_0$ at the
root and allocate the other $A_0$ to later levels by the same rule.
The path coefficient is
\[
 \frac1{\sqrt{A_0}}
 \left(1+\frac{c}{(1-c^{2/3})^{3/2}}\right).
\]
For $\gamma=9/10$ and $r=16$, $c=1/8$ and $c^{2/3}=1/4$.
The root/tail shape in~\eqref{eq:main-schedule} is exactly this shape,
with strict inward shift from denominator $80$ to $81$.
The boundary coefficient is $76.4327736858\ldots$; it is an infimum at
an excluded boundary, not the displayed certified coefficient.

\begin{thebibliography}{99}
\raggedright
\bibitem{upstream} OpenAI, \emph{The Exponential Time Hypothesis is False},
October 25, 2026. Pinned preprint directory:
\sourceurl{https://github.com/openai/math/tree/7ef5c0cc5ca4b0e072b7a050fca09ab24bc6fa3a/preprints/The-Exponential-Time-Hypothesis-is-False-October-25-2026}.
\bibitem{contraction} Same source, \texttt{build/quantitative-contraction.tex},
Git blob \texttt{51ba54efc29f51bbc0426f848003030da484f354}; especially
\texttt{lem:tensor-packets}, \texttt{lem:tapered-hermite},
\texttt{eq:tensor-recursion-identity}, and \texttt{subsec:tensor-children}.
\bibitem{reduction} Same source, \texttt{build/quantitative-reduction.tex},
Git blob \texttt{80389ac39d987133b71b2ddc2593e6fcdd99871a}.
\bibitem{repoindex} Vladimir Reshetnikov, ProveIt ETH project index,
placement commit \texttt{2b925d8800cda68f5723b65dc1fcdaf3d8170359}.
\sourceurl{https://github.com/VladimirReshetnikov/ProveIt/blob/2b925d8800cda68f5723b65dc1fcdaf3d8170359/Computability/SubexponentialSAT/README.md}.
\bibitem{cauchypacket} \emph{Factorial Cauchy modes for the subexponential
3-SAT algorithm}, October 26, 2026, \texttt{Research/cauchy\_mode\_improvement}
in the pinned ProveIt project. Arrival archive:
\texttt{cauchy\_mode\_improvement\_2026-10-26.zip}.
\bibitem{contractionpacket} \emph{Quantitative improvements to the
exact-counting contraction algorithm}, October 26, 2026,
\texttt{Research/eth\_contraction\_improvements}; arrival archive
\texttt{eth\_contraction\_improvements.zip}.
\bibitem{precisionpacket} \emph{Sharper precision for subexponential 3SAT},
October 26, 2026, \texttt{Research/3sat\_precision\_refinement}; arrival
archive \texttt{3sat\_precision\_refinement\_2026-10-26.zip}.
\bibitem{densitypacket} \emph{Improving the clause-density dependence of
exact k-SAT counting}, October 26, 2026,
\texttt{Research/sat\_density\_improvement}; arrival archive
\texttt{sat\_density\_improvement.zip}.
\bibitem{densityrefinement} \emph{A density-sensitive refinement of exact
\#k-SAT counting}, October 26, 2026,
\texttt{Research/eth\_density\_refinement}; arrival archive
\texttt{eth\_density\_refinement.zip}.
\bibitem{scalarpacket} \emph{Refining the October 25, 2026 subexponential
SAT algorithm}, October 26, 2026, \texttt{Research/3sat\_refinement};
arrival archive \texttt{3sat\_refinement\_package.zip}.
\bibitem{barrierpacket} \emph{Asymptotic improvement analysis for the
October 2026 3 SAT preprint}, October 26--27, 2026,
\texttt{Research/3sat\_asymptotic\_analysis}; arrival archive
\texttt{3sat\_asymptotic\_analysis\_20261027.zip}.
\bibitem{claudepaper} Claude (Anthropic), \emph{A $2^{o(n)}$ time algorithm
for 3-SAT}, October 26, 2026, 3:13 PM Pacific, 49 pages. Delivered PDF
at \texttt{Research/claude-2o-n-3sat/eth.pdf} in the pinned project.
\bibitem{leanscope} OpenAI, scope page for the upstream Lean development:
\sourceurl{https://github.com/openai/math/blob/7ef5c0cc5ca4b0e072b7a050fca09ab24bc6fa3a/lean/docs/378.md}.
\end{thebibliography}
\end{document}
